
%===================================================================================================%

\frame{\frametitle{The issue}
Eventually you will come across data from many different
people/groups. Possible problems:
\begin{itemize}
\item Different conventions
\item Different data formats
\item Duplicate entries
\item Damaged/lost/unlabelled files
\end{itemize}
\vspace{5mm}
\centering
% FAIR paper
We are in an era where significant effort is put into
ensuring \\ data can be \textbf{readily used and reproduced}.
}

\frame{\frametitle{Practical goals for your own work}
What will save you a lot of pain:
\begin{itemize}
\item You know where your data is
\item You can read your old data
\item Sufficient metadata to know what your data means
\item You can tell if your data got damaged 
\item Your data is homogeneous and/or well organized
\end{itemize}
}


\frame{\frametitle{Toward these goals}
\begin{enumerate}
\item Bring input data into common format
% Example: Graduate student asks me to analyze data,
% format different between data sets
\begin{itemize}
\item Makes it easier to write/debug scripts
\item Helps next person using your work
\item Opportunity for cross-checks
\end{itemize}
\item Stay consistent with conventions
\item Try to use already-existing data structures (e.g.
\texttt{yaml} or \texttt{csv})
\begin{itemize}
\item Everyone can leverage already-existing tools
\item Lots of information about format structure 
\end{itemize}
\item Label your columns
\item Add \texttt{README.md} for supplemental information
\end{enumerate}
}


